\section{A sine-normalized cyclotomic family}\label{sec:cyclotomic}
Algebraic approximations offer another route: choose an algebraic number
$\alpha$ close to $\theta$ and evaluate its integer minimal polynomial
at $\theta$. This evaluation contains a factor $\theta-\alpha$, but
also a factor for every other conjugate of $\alpha$. A good approximation
at one real embedding can therefore be defeated by large conjugates.
The roots-of-unity family below exhibits this explicitly, including
the removal of common factors from the polynomial coefficients.

The topological-signature motivation came from the cubic cotangent
formula recorded by March\'e and Masbaum~\cite[Corollary 3.2]{MarcheMasbaum},
who credit its simplification to Murakami. That signature formula is not
an even-power integrality theorem and is not used below.
The following algebraic and analytic claims are derived directly for
one specified normalization.

For an odd prime $p$, set $h=\pi/(2p)$, $w=e^{2\pi i/(4p)}$, and
\begin{equation}\label{eq:alpha}
 \alpha_p=\sin^4h\sum_{\substack{1\leq n<p\\n\ {
m odd}}}
             \chi_4(n)\frac{\cos^3(nh)}{\sin^4(nh)}.
\end{equation}
Let $F_p=\Q(w+w^{-1})$, of degree $d=p-1$, and
$M_p(X)=\operatorname{Norm}_{F_p/\Q}(X-\alpha_p)$.
Here the norm polynomial is the product of $X-\sigma(\alpha_p)$
over all embeddings $\sigma:F_p\hookrightarrow\R$.

\begin{theorem}[One cyclotomic norm obstruction; S06]\label{thm:norm}
The number $8\alpha_p$ is integral, $\alpha_p$ has degree $p-1$,
and its primitive minimal polynomial $P_p$ satisfies
\[
 \log|P_p(\theta)|=4(p-1)\log p+O(p).
\]
In particular these primitive evaluations grow along the odd primes.
\end{theorem}
\begin{proof}
The identity
\[
 8\alpha_p=\sum_{\substack{n<p\\n\ {
m odd}}}\chi_4(n)
 \left(\frac{w-w^{-1}}{w^n-w^{-n}}\right)^4(w^n+w^{-n})^3
\]
proves integrality. Indeed $z=w^2$ is a primitive $2p$th root and
$1-z^n$ is a cyclotomic unit for odd $0<n<p$, since
$\Phi_{2p}(1)=1$. Polynomial division of
$\Phi_{2p}(x)-\Phi_{2p}(1)$ by $x-1$ proves the unit assertion.
The displayed expression is real and belongs to $F_p$.

The principal approximation has the signed expansion
\begin{equation}\label{eq:alpha-error}
 \alpha_p=\theta-
 \frac{\pi^2(4\theta+5\beta(2))}{24p^2}+O(p^{-4}).
\end{equation}
To prove it, write
$\cos^3x/\sin^4x=x^{-4}-(5/6)x^{-2}+R(x)$, where $R$ is smooth
on $[0,\pi/2]$. Alternating samples of $R$ have uniformly bounded
sum by its total variation, and the alternating tails of $n^{-4}$
and $n^{-2}$ are bounded by $p^{-4}$ and $p^{-2}$. Multiply by
$\sin^4h$ and use $(\sin h/h)^4=1-(2/3)h^2+O(h^4)$.
The coefficient of $p^{-2}$ is strictly positive.

Represent the real embeddings by odd $1\leq k<2p$, $k\ne p$.
The conjugate $\alpha_{p,k}$ is given by~\eqref{eq:alpha} with each
angle multiplied by $k$. Let
\[
 f(x)=\frac{\cos^3x}{\sin^4x},\qquad
 C=\sum_{n\geq1,\ n\ {
m odd}}n^{-4}=\frac{\pi^4}{96},
 \qquad\delta=2-C>0.
\]
The function $x^4f(x)$ strictly decreases on $(0,\pi/2)$.
The sign of its logarithmic derivative reduces to positivity of
$H(x)=x(3+\cos^2x)-4\sin x\cos x$; indeed $H(0)=0$ and
$H'(x)=7\sin^2x-2x\sin x\cos x\geq5\sin^2x$.
Thus $f(mh)\leq f(h)/m^4$ for odd $1\leq m<p$.

Multiplication by $k$ permutes the pairs of odd nonzero residues
$\{n,2p-n\}$ modulo $2p$. The inner conjugate sum is therefore
$\sum_m\eta_m f(mh)$, where each $\eta_m=\pm1$ and each odd
$0<m<p$ occurs once. Its first term exceeds the sum of the absolute
values of the rest. Uniformly in $k$,
\begin{equation}\label{eq:conjugate-bounds}
 \delta\cos^3h\left(\frac{\sin(kh)}{\sin h}\right)^4
 \leq|\alpha_{p,k}|\leq
 C\cos^3h\left(\frac{\sin(kh)}{\sin h}\right)^4.
\end{equation}
For $p\geq5$ and $k\notin\{1,2p-1\}$, the sine ratio is at least
$\sin(3h)/\sin h\geq2$, so these conjugates have absolute value
greater than $10$ (use $\delta>0.98$, $\cos^3h\geq3\sqrt3/8$).
The principal conjugate is in $(0,C)$ and the $2p-1$ conjugate is
its negative. Hence it is the unique conjugate in $(0,2)$ and its
stabilizer is trivial. The degree is $d$. For $p=3$, directly
$\alpha_3=3\sqrt3/8$ has degree two.

The odd sine product is
\[
 \prod_{\substack{1\leq k<2p\\k\ {
m odd},\ k\ne p}}\sin(kh)
       =2^{1-p}=2^{-d}.
\]
It follows by evaluating $x^p+1$ at $x=1$ in its root factorization
and removing the sine-one factor. Therefore~\eqref{eq:conjugate-bounds}
implies
\[
 \left(\frac{\delta\cos^3h}{16\sin^4h}\right)^d
 \leq\prod_k|\alpha_{p,k}|\leq
 \left(\frac{C\cos^3h}{16\sin^4h}\right)^d.
\]
For nonprincipal conjugates, replacing $|\alpha_{p,k}|$ by
$|\theta-\alpha_{p,k}|$ changes the product logarithm by $O(d)$,
since $0<\theta<1$ and their magnitudes exceed $10$.
For the two remaining conjugates,~\eqref{eq:alpha-error} gives a
total change $-2\log p+O(1)$. Since $\sin h\sim\pi/(2p)$,
$\log|M_p(\theta)|=4d\log p+O(p)$.

Finally $J_p(X)=\operatorname{Norm}(8X-8\alpha_p)=8^dM_p(X)$
is integral, and its content $g_p$ divides its leading coefficient
$8^d$. Thus $P_p=J_p/g_p$ is the actual primitive minimal polynomial
and $|M_p(\theta)|\leq|P_p(\theta)|\leq8^d|M_p(\theta)|$.
This proves the theorem after all content is removed.
\end{proof}

The norm polynomial is evaluated at the real number $\theta$; no
assertion that $\theta-\alpha_p\in F_p$ is made without first
assuming $\theta$ rational. Under $\theta=a/b$, the integer test is
$b^dP_p(a/b)\in\Z\setminus\{0\}$. Its degree cost cannot be omitted.
Neither other normalizations nor other nonlinear norms are excluded.
